\subsection{Nonlinear mechanical dynamics}
The direct rate-product construction is especially natural for mechanical
systems. A robot with joint positions $q$, velocities $w=\dot q$, joint
torques $\tau_m$, and an external contact wrench $\lambda$ obeys, on a
smooth interval without impacts, 
\begin{equation}
M(q)\ddot q+C(q,\dot q)\dot q+G(q) =\tau_m+J_c(q)^{\mathsf{T}}\lambda.
\label{eq:mechanical}
\end{equation}
The inertia matrix $M(q)$ couples joint accelerations, the term $C(q,w)w$
describes velocity-dependent forces, $G(q)$ accounts for gravity, and the
contact Jacobian $J_c(q)$ maps the external wrench into generalized joint
forces. A command $u$ produces torque $\tau_m=B(q)u$, where the actuation
matrix $B(q)$ need not be square; for direct torque actuation $B=I$. Write $%
h(q,w)=C(q,w)w+G(q)$. The force balance is then 
\begin{equation}
M(q)\dot w+h(q,w)-B(q)u-J_c(q)^{\mathsf{T}}\lambda=0.
\label{eq:mechanical_balance}
\end{equation}

To express this balance in terms of context rates, retain $q$ and $w$ and
add the integrated command $s_u$, integrated contact wrench $s_\lambda$, and
clock coordinate $r_0$ defined by $\dot s_u=u$, $\dot s_\lambda=\lambda$,
and $\dot r_0=1$. For the resulting context 
\begin{equation}
z=\func{col}(q,w,s_u,s_\lambda,r_0),  \label{eq:mechanical_context}
\end{equation}
the observed or estimated rates are 
\begin{equation}
\dot z=\func{col}(w,\dot w,u,\lambda,1).  \label{eq:mechanical_rates}
\end{equation}
We now collect the coefficients of these rates, in the same order, into the
matrix-valued function 
\begin{equation}
H(z)=%
\begin{bmatrix}
0 & M(q) & -B(q) & -J_c(q)^{\mathsf{T}} & h(q,w)%
\end{bmatrix}%
.  \label{eq:mechanical_coefficients}
\end{equation}
The original equations of motion are therefore exactly 
\begin{equation}
H(z)\dot z=0.  \label{eq:mechanical_product}
\end{equation}
The symbol $H$ denotes the coefficient matrix assembled from the physical
equations; it is not introduced as the derivative of an unknown function. In
general it need not equal a Jacobian $D_zf$: such an equality would require
the additional integrability conditions on its entries. Nevertheless, the
rate-product equation in (\ref{eq:mechanical_product}) supplies a
differential compatibility relation directly, without differentiating the
force balance in (\ref{eq:mechanical_balance}). For the learning agent, the
mechanics and their coefficient blocks may be unknown; the measured motion
and commanded inputs provide the variations from which a compatible relation
can be estimated. Contact forces must themselves be measured or estimated if
they are to enter the chosen context. The additional Acrobot example below uses a prescribed mechanical family.

